% Source-grounded companion to docs/RAG检索与跨文本证据设计.md.
\section{跨文本举措定位与系统检索设计}

工作台把资料整理、跨文本举措定位和原文检索连到前置模式判定、六维观察、市场收益、版本变化与专项调查。用户选择一项举措，可以对照不同材料中的提议、实际决定、实施和结果，回到原文及相邻上下文；也可以把已保存解析材料导出为Markdown或JSON，供后续检索使用。原文检索先限定项目、版本、来源和时点，再复用现有中文分词与BM25排序。工程量核算、月度运营计算和神经网络造价预测继续使用各自审核输入，检索文字不直接写入金额或当前事实。

\subsection{从解析材料到可回读证据}

导出以既有解析快照为输入，按文档标识、哈希版本、原文定位组织正文、段落、页面、文本行和表格行。优先读取原始解析结果，不将界面表格预览当作全表；原解析只保存预览时，明确登记已导出行数和截取范围。扫描件没有可读文本时保留解析问题。Markdown保存可阅读原文及引用，JSON保存相同材料身份、原文片段和审核字段，上传原件保持不变。

\begin{longtable}{L{25mm}L{50mm}L{80mm}}
\toprule
索引对象 & 保留内容 & 在研判中的作用\\
\midrule
\endhead
文档与来源 & 文档标识、文件名、哈希、来源类别、来源登记与解析覆盖 & 对应原件身份和当前版本；区分公开原件、内部材料、AI草稿、模拟示例和未知来源。\\
原文片段 & 原文定位、完整文字、表格行与字段值 & 按页、段、行或表格回读；未被案例引用的上下文也参与检索。\\
举措引用 & 举措编号、卡号与修订、字段、阶段、对象、审核状态和逐字引文 & 同一举措跨文本对照，分别保留提议、决定、实施、结果、条件与排除条件。\\
文字提及 & 命中举措名称、字符位置及未审核标记 & 提供查找线索，阶段与对象由原记录另行确认。\\
未绑定引用 & 原件缺失、版本不符、定位缺失或引文不符 & 显示需补齐的证据，不绑定到错误材料。\\
\bottomrule
\end{longtable}

已登记举措引用同时核对同文档、同哈希、同定位和连续逐字引文。正文只出现“二次供水”等名称时，显示为未复核的文字提及；其出现次数不表示改造次数或独立印证数量。文本解析生成的行文字和\texttt{Line/Text}表格镜像，在逐字一致且没有额外独立引用时合并检索命中，并保留别名定位；导出仍保存原解析定位，独立文件和实际表格记录保持各自身份。

导出沿用原审核状态，用途登记为检索准备，不确认训练权利。当前导出契约没有许可授权字段，后续训练或对外分发前需登记允许用途、授权范围和责任人；训练用途默认未确认。

\subsection{范围筛选、词项排序与系统入口}

现有原文检索接口接受问题、举措编号、对象标识、原文登记阶段、来源类别、是否浏览AI材料、审核状态、资料截止日期和最多5条结果。服务端先检查项目所有者，再核对原件哈希、解析源哈希及项目内来源绑定。举措、对象、阶段和审核条件必须在同一条实际引用中同时成立，不能把同段不同引用的条件拼成一个新对象。

默认检索排除AI研究草稿与模拟示例，显式开启后分别标明类别。资料截止日期采用核实后的正文发布日期，上传或采集日不替代正文时点，晚于截止日、日期未核或超过登记有效期的来源保留排除原因。没有日期筛选的原文浏览可以展示日期未知的材料，并显示缺项。

排序直接复用现有\texttt{jieba.cut}与\texttt{BM25Okapi}。当前只增加三组明确措辞桥接：“楼栋/楼幢”补充“幢”、“户数/住户”补充“户”、“排除”补充“不纳入”；原问题、实际排序问题和扩展原因同时保留。BM25分数随词频、语料和片段长度改变，也可能为负数，表示本次词项相关程度，不表示事实可信率、工程效果或举措适用率。当前没有向量库或语义检索效果评价。

全部符合范围的已保存片段参加候选构造，包括上下文和超出界面预览的表格行。大资料集按文档交错分批，每批最多1000片段及约1MB文本，各批前5候选汇总后用同一BM25再次排序，批次分数不直接比较。结果显示原文、来源日期、版本、定位、别名定位、举措、对象、阶段、审核状态与命中原因；范围不符的材料另列排除原因。整个检索过程不写入项目事实，正文中的角色、工具或外发指令作为材料文字保存，不改变运行任务与权限。

系统查询围绕五类子问题组织：居民需求寻找问卷对象与采集期；举措追踪分别寻找决定、公示、施工和验收；规则适用核对有效条文与当前条件；费用比较核对工程范围、计价阶段、面积口径与价格年；运营评价核对月份、独立使用人数、人次、购买服务、补贴及直接支付。每个子问题的举措、对象、目标字段、时点和所需阶段可组织成业务查询规划，再映射到现有接口。任务自动拆分与目标字段路由属于后续接入。

自动研判继续走已有已审核知识入口：案例与规则按所有者、研判时点、真实/模拟范围、公开/内部模式和条件匹配，运行时冻结本次知识版本，DSH Agent读取命中依据及现实项目材料，程序逐项核对举措、命中卡、已匹配条件和连续逐字引用。最终决策同时要求现实来源依据，硬约束与关键缺项保持原校验流程。新增原文浏览接口尚未自动接入DSH原文问答生成；后续接入复用同一运行时与引用验证，不新增调度器。

\subsection{求是村范围例与调查落点}

现有公开原件《浙大求是社区推进二次供水改造》日期为2023年8月18日。\texttt{line:22}记载拟改29、30、31、48幢，共252户，并记载留博楼因当时泵房位置问题未纳入；\texttt{line:24}记录方案与排除原因公示、供水合同签订及后续拟进场施工。查询“二次供水改造涉及哪些楼栋和户数，哪些楼栋被排除”，能够回读相关片段，并区分公开原文与尚未审核的提取字段。252户属于该报道中的四幢工程范围，阶段为报道时点的进场准备，不据此登记2026年已经改造完成。

若用户问当前完成情况，调查任务应对应四幢的施工、验收、接管与运行记录，留博楼单列核查；任务列明对象、需取得文件、待补字段和影响方案。同名AI研究报告保持独立材料身份，作为回找原件的线索；重复转述不增加独立依据。真实接口核验已检查公开/AI分离、同版本定位、截止日期排除和项目输入不变，事实写入为0，未调用语言模型。

\subsection{后续效果评价}

当前验证为工程回归与真实本地接口示例。后续按项目和研判时点建立人工标注问题集，覆盖同一举措跨文本、多对象同段、计划与实施差异、政策正式稿与征求稿、转载、金额阶段及资料缺失。标注每个问题应返回的目标片段、有效引用、对象、阶段、日期和必须调查事项；生成答案不作为标注真值。

\begin{longtable}{L{29mm}L{126mm}}
\toprule
评价指标 & 计算对象与作用\\
\midrule
\endhead
Recall@5 & 前5条命中覆盖的有效目标片段数除以全部有效目标片段数，检查措辞扩展、分批与跨文本召回。\\
引用准确率 & 同时满足原文、版本、定位、对象和阶段的引用条数除以全部输出引用条数。\\
范围与日期误用率 & 错用楼幢、阶段、截止日期之后或失效资料的查询比例，分别报告。\\
AI/模拟误用率 & 将AI草稿或模拟数据表述为真实依据的查询比例。\\
调查覆盖率 & 资料不足问题中正确列出必须补齐目标字段的查询比例。\\
\bottomrule
\end{longtable}

上述指标尚无人工标注统计结果。先固定材料、问题与标注，再比较现有BM25、措辞扩展和分批策略；发现同义表达的召回缺口后，才在同一权限、版本和索引契约下评估语义召回。实现依据为\texttt{packages/experience}中的解析切片、导出与检索模块及现有知识研判入口。设计细节见docs目录的《RAG检索与跨文本证据设计》Markdown文档。

真实接口证据保存于\path{research/test_results/rag-workspace-live.json}。
